%-------------------------------------------------------------------------------
%   CHAPTER 3 - СУДАЛГААНЫ АРГА ЗҮЙ
%-------------------------------------------------------------------------------

\section{Судалгааны арга зүй}
\label{sec:methodology}

Энэ бүлэгт Egüü системийн шаардлагыг тодорхойлж, системийн шинжилгээ, зохиомж, архитектурын шийдлүүдийг дэлгэрэнгүй танилцуулна. Мөн хэрэглэгчийн интерфэйс (UI/UX) загвар, өгөгдлийн сангийн зохиомжийг авч үзнэ.

\subsection{Системийн зохиомжийн шийдлүүд}
\label{subsec:design-decisions}

Энэ хэсэгт Egüü системийг хэрхэн зохион бүтээх, ямар технологи ашиглах, яагаад тэр технологиудыг сонгосон талаар тодорхой тайлбарлана.

\subsubsection{Гадаад зохиомж (External Design)}

\textbf{Хэрэглэгчийн харагдах хэсэг:} Egüü систем нь хэрэглэгчид tab-based навигацитай 5 үндсэн дэлгэц (Нүүр, Хайлт, Дуртай, Профайл, Камер) өгч, flashcard-уудыг категориор (эгшиг, гийгүүлэгч, үг) ангилж, хайлтын функцээр хялбар олох боломжийг олгодог бөгөөд энэ нь хэрэглэгчид ойлгомжтой, интуитив туршлага өгдөг.

\paragraph{Гадаад зохиомжийн үндсэн шийдлүүд:}

\begin{itemize}
    \item \textbf{Навигаци:} Tab-based навигаци ашигласан нь хэрэглэгч хамгийн түгээмэл хэрэглэгддэг дэлгэцүүд рүү нэг товшилтоор шилжих боломжтой болгодог
    
    \item \textbf{Категоризаци:} Flashcard-уудыг категориор ангилсан нь хэрэглэгч өөрийн сонирхсон хэсгээс эхлэн суралцах боломжтой
    
    \item \textbf{Хайлт:} Фонетик, англи орчуулга, тодорхойлолтоор хайх боломжтой нь хэрэглэгч хурдан олох боломжтой
    
    \item \textbf{Gamification элементүүд:} Streak counter, daily word, progress bar зэрэг нь хэрэглэгчийн мотивацийг нэмэгдүүлдэг
\end{itemize}

\subsubsection{Дотоод зохиомж (Internal Design)}

\textbf{Системийн дотоод бүтэц:} Систем нь client-server архитектураар зохион байгуулагдсан бөгөөд React Native клиент нь Firebase Backend (Authentication, Firestore, Storage) болон TensorFlow.js OCR модультай REST API болон SDK-аар харилцаж, бизнес логик нь клиент талд, өгөгдлийн хадгалалт нь Firebase үүлэн дээр байрладаг.

\paragraph{Дотоод зохиомжийн үндсэн бүрэлдэхүүн хэсгүүд:}

\begin{enumerate}
    \item \textbf{Client Layer (React Native):}
    \begin{itemize}
        \item Presentation Layer: UI компонентууд (FlashcardItem, StreakCounter, DailyWordCard)
        \item Business Logic Layer: Spaced repetition алгоритм, streak tracker, validators
        \item Data Access Layer: Firebase SDK, TensorFlow.js wrapper
    \end{itemize}
    
    \item \textbf{Backend Layer (Firebase):}
    \begin{itemize}
        \item Authentication Service: Хэрэглэгчийн баталгаажуулалт
        \item Database Service: Firestore NoSQL өгөгдлийн сан
        \item Storage Service: Зураг, файл хадгалах
        \item Security Rules: Өгөгдлийн хандалтын хяналт
    \end{itemize}
    
    \item \textbf{ML Layer (TensorFlow.js):}
    \begin{itemize}
        \item OCR Model: Монгол бичиг таних загвар
        \item Image Preprocessing: Зураг боловсруулах
        \item Prediction Engine: Таних үйлдэл гүйцэтгэх
    \end{itemize}
\end{enumerate}

\subsubsection{Өгөгдлийн сангийн сонголт}

\textbf{Cloud Firestore NoSQL сонголт:} Cloud Firestore NoSQL өгөгдлийн санг сонгосон учир нь real-time синхрончлол (хэрэглэгчийн өөрчлөлтүүд шууд бүх төхөөрөмжид харагдах), offline persistence (интернэтгүй үед ч ажиллах), автомат scalability (хэрэглэгчийн тоо нэмэгдэхэд автоматаар өргөжих), мөн document-based бүтэц нь Монгол бичгийн flashcard өгөгдлийг (mongolBichig, phonetic, english, definition, imageUrl гэх мэт олон талбартай) уян хатан хадгалахад тохиромжтой.

\paragraph{Firestore vs бусад өгөгдлийн сангийн харьцуулалт:}

\begin{table}[H]
\centering
\caption{Өгөгдлийн сангийн харьцуулалт}
\label{tab:database-comparison}
\begin{tabular}{|l|c|c|c|}
\hline
\textbf{Шинж чанар} & \textbf{Firestore} & \textbf{MongoDB} & \textbf{PostgreSQL} \\
\hline
Real-time sync & +++ & ++ & - \\
\hline
Offline support & +++ & + & - \\
\hline
Auto-scaling & +++ & ++ & + \\
\hline
Setup хялбар байдал & +++ & ++ & + \\
\hline
Mobile SDK & +++ & + & - \\
\hline
Үнэ (жижиг төсөл) & +++ & ++ & ++ \\
\hline
Нарийн query & + & +++ & +++ \\
\hline
\end{tabular}
\end{table}

\subsubsection{Технологийн стек сонголт}

\textbf{Ашигласан технологиуд ба шалтгаан:} React Native (cross-platform хөгжүүлэлт, нэг кодоор iOS/Android дээр ажиллах), Firebase (backend-as-a-service, сервер удирдах шаардлагагүй), TensorFlow.js (client-side ML, privacy хадгалах), NativeWind (Tailwind CSS for React Native, хурдан UI хөгжүүлэлт), react-i18next (олон хэлний дэмжлэг, Монгол/Англи хэл солих) зэрэг орчин үеийн технологиудыг сонгосон нь хөгжүүлэлтийн хурд, найдвартай байдал, хэрэглэгчийн туршлагыг сайжруулж, засвар үйлчилгээний зардлыг бууруулдаг.

\paragraph{Технологи бүрийн сонголтын шалтгаан:}

\begin{table}[H]
\centering
\caption{Технологийн сонголтын шалтгаан}
\label{tab:tech-stack-reasons}
\begin{tabular}{|l|p{5cm}|p{6cm}|}
\hline
\textbf{Технологи} & \textbf{Давуу тал} & \textbf{Egüü төсөлд яагаад тохиромжтой} \\
\hline
React Native & Cross-platform, Hot Reload, Том коммьюнити & Нэг удаа бичээд iOS/Android дээр ажиллуулах, хөгжүүлэлтийн цаг 50\% хэмнэнэ \\
\hline
Firebase & BaaS, Real-time, Auto-scaling & Сервер удирдах шаардлагагүй, бага төсөвт тохиромжтой \\
\hline
TensorFlow.js & Client-side ML, Privacy & Хэрэглэгчийн өгөгдөл сервер рүү явахгүй, offline ажиллана \\
\hline
NativeWind & Tailwind CSS, Хурдан UI & Utility-first CSS, responsive дизайн хялбар \\
\hline
i18next & Олон хэлний дэмжлэг & Монгол/Англи хэл солих, цаашид бусад хэл нэмэх хялбар \\
\hline
Expo & Managed workflow, EAS Build & Xcode/Android Studio суулгах шаардлагагүй, хурдан туршилт \\
\hline
\end{tabular}
\end{table}



\subsection{Системийн шаардлага}
\label{subsec:requirements}

\subsubsection{Функциональ шаардлага}

Функциональ шаардлага нь системийн хийх ёстой үндсэн үйлдлүүдийг тодорхойлдог. Egüü системд дараах функциональ шаардлагууд тавигдана:

\paragraph{Хэрэглэгчийн удирдлага:}
\begin{itemize}
    \item FR-1: Хэрэглэгч Email/Password-оор бүртгүүлэх, нэвтрэх боломжтой байх
    \item FR-2: Хэрэглэгч Google account-аар нэвтрэх боломжтой байх
    \item FR-3: Хэрэглэгч профайл мэдээллээ засах боломжтой байх
    \item FR-4: Хэрэглэгч нууц үгээ сэргээх боломжтой байх
\end{itemize}

\paragraph{Flashcard сургалтын систем:}
\begin{itemize}
    \item FR-5: Систем нь Монгол бичгийн үсгүүдийг flashcard хэлбэрээр харуулах
    \item FR-6: Flashcard нь зураг, фонетик, англи орчуулга, тодорхойлолт агуулах
    \item FR-7: Хэрэглэгч flashcard-ыг эргүүлж хариултыг харах боломжтой байх
    \item FR-8: Систем нь spaced repetition алгоритмаар дараагийн давталтын хугацааг тооцоолох
    \item FR-9: Хэрэглэгч flashcard-д "амархан", "дунд", "хэцүү" гэх мэтээр үнэлгээ өгөх боломжтой байх
\end{itemize}

\paragraph{Ангилал ба хайлт:}
\begin{itemize}
    \item FR-10: Flashcard-уудыг категориор (эгшиг, гийгүүлэгч, үг гэх мэт) ангилах
    \item FR-11: Хэрэглэгч flashcard хайх боломжтой байх (фонетик, англи орчуулгаар)
    \item FR-12: Категори бүрээр шүүх боломжтой байх
\end{itemize}

\paragraph{Дуртай жагсаалт:}
\begin{itemize}
    \item FR-13: Хэрэглэгч flashcard-ыг дуртай жагсаалтдаа нэмэх боломжтой байх
    \item FR-14: Дуртай flashcard-уудыг тусдаа харах боломжтой байх
    \item FR-15: Дуртай жагсаалтаас хасах боломжтой байх
\end{itemize}

\paragraph{Өдөр бүрийн үг (Daily Word):}
\begin{itemize}
    \item FR-16: Систем нь өдөр бүр шинэ үг санал болгох
    \item FR-17: Хэрэглэгч өдрийн үгийг дэлгэрэнгүй үзэх боломжтой байх
    \item FR-18: Өдрийн үгийн түүхийг хадгалах
\end{itemize}

\paragraph{Streak (Дараалсан өдрүүд) tracking:}
\begin{itemize}
    \item FR-19: Систем нь хэрэглэгчийн дараалсан суралцсан өдрүүдийг тооцоолох
    \item FR-20: Streak тасарсан тохиолдолд мэдэгдэл өгөх
    \item FR-21: Streak-ийн статистикийг харуулах
\end{itemize}

\paragraph{OCR (Гар бичлэг таних):}
\begin{itemize}
    \item FR-22: Хэрэглэгч камераар Монгол бичгийн үсэг зурж таниулах боломжтой байх
    \item FR-23: Систем нь TensorFlow.js ашиглан үсгийг таних
    \item FR-24: Таниулсан үсгийн үр дүнг харуулах
    \item FR-25: Таниулсан үсгийн түүхийг хадгалах
\end{itemize}

\paragraph{Тохиргоо ба хэл солих:}
\begin{itemize}
    \item FR-26: Хэрэглэгч интерфэйсийн хэлийг солих боломжтой байх (Монгол, Англи)
    \item FR-27: Dark/Light theme солих боломжтой байх
    \item FR-28: Мэдэгдлийн тохиргоог өөрчлөх боломжтой байх
\end{itemize}

\subsubsection{Функциональ бус шаардлага}

Функциональ бус шаардлага нь системийн чанар, гүйцэтгэл, аюулгүй байдал зэрэг шинж чанаруудыг тодорхойлдог.

\paragraph{Гүйцэтгэл (Performance):}
\begin{itemize}
    \item NFR-1: Апп нь 3 секундээс бага хугацаанд ачаалагдах ёстой
    \item NFR-2: Flashcard эргүүлэх үйлдэл нь 100ms-ээс бага хугацаанд ажиллах
    \item NFR-3: Хайлтын үр дүн нь 500ms-ээс бага хугацаанд харагдах
    \item NFR-4: OCR таних үйлдэл нь 2 секундээс бага хугацаанд дуусах
\end{itemize}

\paragraph{Хэрэглэхэд хялбар байдал (Usability):}
\begin{itemize}
    \item NFR-5: Интерфэйс нь ойлгомжтой, интуитив байх
    \item NFR-6: Шинэ хэрэглэгч 5 минутын дотор үндсэн функцуудыг ашиглаж сурах
    \item NFR-7: Бүх товч, текст нь тодорхой, уншихад хялбар байх
    \item NFR-8: Алдааны мэдэгдэл нь ойлгомжтой, шийдлийн зөвлөмж өгөх
\end{itemize}

\paragraph{Найдвартай байдал (Reliability):}
\begin{itemize}
    \item NFR-9: Систем нь 99\% uptime хангах
    \item NFR-10: Өгөгдөл алдагдахгүй байх (Firebase backup)
    \item NFR-11: Offline горимд ажиллах боломжтой байх
    \item NFR-12: Алдаа гарсан тохиолдолд хэрэглэгчид ойлгомжтой мэдэгдэл өгөх
\end{itemize}

\paragraph{Аюулгүй байдал (Security):}
\begin{itemize}
    \item NFR-13: Хэрэглэгчийн нууц үг нь hash хийгдсэн байх
    \item NFR-14: Firebase Security Rules ашиглан өгөгдлийн хандалтыг хязгаарлах
    \item NFR-15: HTTPS протокол ашиглах
    \item NFR-16: Хэрэглэгчийн хувийн мэдээлэл гуравдагч талд дамжихгүй байх
\end{itemize}

\paragraph{Өргөтгөх боломж (Scalability):}
\begin{itemize}
    \item NFR-17: Систем нь 10,000+ хэрэглэгчид үйлчлэх боломжтой байх
    \item NFR-18: Flashcard dataset-ийг хялбар нэмэх боломжтой байх
    \item NFR-19: Шинэ хэл нэмэх боломжтой байх (i18n)
    \item NFR-20: Шинэ функц нэмэхэд хялбар байх (модуляр архитектур)
\end{itemize}

\subsection{UML диаграммууд}
\label{subsec:uml-diagrams}

Энэ хэсэгт Egüü системийн үйл ажиллагааг янз бүрийн өнцгөөс харуулсан UML (Unified Modeling Language) диаграммуудыг танилцуулна. Эдгээр диаграммууд нь системийн бүтэц, үйл явц, өгөгдлийн урсгалыг ойлгомжтой байдлаар харуулдаг.

\subsubsection{Use Case диаграмм}

Use Case диаграмм нь системийн үндсэн функцуудыг болон хэрэглэгчийн харилцааг харуулдаг. Egüü системд дараах үндсэн use case-ууд байна:

\begin{figure}[H]
    \centering
    \begin{tikzpicture}[
        actor/.style={stick figure, draw=black, minimum width=1cm, minimum height=2cm},
        usecase/.style={ellipse, draw=black, minimum width=3cm, minimum height=1cm, align=center},
        system/.style={rectangle, draw=black, minimum width=10cm, minimum height=12cm}
    ]
    
    % System boundary
    \node[system] (system) at (5, 0) {};
    \node[above] at (system.north) {\textbf{Egüü System}};
    
    % Actor
    \node[actor] (user) at (-2, 0) {};
    \node[below] at (user.south) {Хэрэглэгч};
    
    % Use cases
    \node[usecase] (register) at (5, 5) {Бүртгүүлэх/\\Нэвтрэх};
    \node[usecase] (study) at (5, 3) {Flashcard\\суралцах};
    \node[usecase] (search) at (5, 1) {Flashcard\\хайх};
    \node[usecase] (favorite) at (5, -1) {Дуртай\\жагсаалт};
    \node[usecase] (daily) at (5, -3) {Өдрийн үг\\үзэх};
    \node[usecase] (ocr) at (5, -5) {Гар бичлэг\\таниулах};
    \node[usecase] (profile) at (9, 0) {Профайл\\удирдах};
    
    % Associations
    \draw (user) -- (register);
    \draw (user) -- (study);
    \draw (user) -- (search);
    \draw (user) -- (favorite);
    \draw (user) -- (daily);
    \draw (user) -- (ocr);
    \draw (user) -- (profile);
    
    \end{tikzpicture}
    \caption{Egüü системийн Use Case диаграмм}
    \label{fig:use-case}
\end{figure}

\paragraph{Use Case диаграммын тайлбар:}

Диаграмм нь системийн 7 үндсэн функцийг харуулж байна: Бүртгүүлэх/Нэвтрэх, Flashcard суралцах, Flashcard хайх, Дуртай жагсаалт, Өдрийн үг үзэх, Гар бичлэг таниулах, Профайл удирдах.

\subsubsection{Sequence Diagram - Хэрэглэгч нэвтрэх процесс}

Sequence диаграмм нь системийн объектууд хоорондын харилцааг цаг хугацааны дагуу харуулдаг. Доорх диаграмм нь хэрэглэгч нэвтрэх процессыг алхам алхмаар харуулж байна.

\begin{figure}[H]
    \centering
    \begin{tikzpicture}[
        node distance=3cm,
        box/.style={rectangle, draw=black, minimum width=2cm, minimum height=0.8cm, align=center},
        message/.style={->, >=stealth, thick}
    ]
    
    % Actors/Objects
    \node[box] (user) {Хэрэглэгч};
    \node[box, right=of user] (ui) {UI\\(Login Screen)};
    \node[box, right=of ui] (auth) {Auth\\Service};
    \node[box, right=of auth] (firebase) {Firebase\\Auth};
    \node[box, right=of firebase] (firestore) {Firestore};
    
    % Lifelines
    \draw[dashed] (user.south) -- ++(0,-10);
    \draw[dashed] (ui.south) -- ++(0,-10);
    \draw[dashed] (auth.south) -- ++(0,-10);
    \draw[dashed] (firebase.south) -- ++(0,-10);
    \draw[dashed] (firestore.south) -- ++(0,-10);
    
    % Messages
    \draw[message] ([yshift=-1cm]user.south) -- node[above, font=\tiny] {1: Email, Password оруулах} ([yshift=-1cm]ui.south);
    
    \draw[message] ([yshift=-2cm]ui.south) -- node[above, font=\tiny] {2: signIn(email, password)} ([yshift=-2cm]auth.south);
    
    \draw[message] ([yshift=-3cm]auth.south) -- node[above, font=\tiny] {3: signInWithEmailAndPassword()} ([yshift=-3cm]firebase.south);
    
    \draw[message] ([yshift=-4cm]firebase.south) -- node[below, font=\tiny] {4: User token} ([yshift=-4cm]auth.south);
    
    \draw[message] ([yshift=-5cm]auth.south) -- node[above, font=\tiny] {5: getUserData(userId)} ([yshift=-5cm]firestore.south);
    
    \draw[message] ([yshift=-6cm]firestore.south) -- node[below, font=\tiny] {6: User data} ([yshift=-6cm]auth.south);
    
    \draw[message] ([yshift=-7cm]auth.south) -- node[below, font=\tiny] {7: Login success} ([yshift=-7cm]ui.south);
    
    \draw[message] ([yshift=-8cm]ui.south) -- node[below, font=\tiny] {8: Нүүр дэлгэц рүү шилжих} ([yshift=-8cm]user.south);
    
    \end{tikzpicture}
    \caption{Хэрэглэгч нэвтрэх процессын Sequence диаграмм}
    \label{fig:sequence-login}
\end{figure}

\paragraph{Sequence диаграммын тайлбар:}

Диаграмм нь нэвтрэх процессыг 8 алхамд хуваан харуулж байна: (1) Email/Password оруулах, (2) Auth Service руу хүсэлт, (3) Firebase баталгаажуулалт, (4) Token буцаах, (5-6) Хэрэглэгчийн өгөгдөл татах, (7-8) Нүүр дэлгэц рүү шилжих. Firebase Authentication ашигласнаар аюулгүй токен-д суурилсан баталгаажуулалт хэрэгждэг.

\subsubsection{Activity Diagram - Flashcard суралцах процесс}

Activity диаграмм нь системийн үйл явцын урсгалыг (workflow) харуулдаг. Энэ нь алхам алхмаар юу болж байгааг, ямар шийдвэр гаргаж байгааг харуулна. Доорх диаграмм нь хэрэглэгч flashcard суралцах процессыг харуулж байна.

\begin{figure}[H]
    \centering
    \begin{tikzpicture}[
        node distance=1.5cm,
        start/.style={ellipse, draw=black, fill=green!30, minimum width=2cm, minimum height=0.8cm, align=center},
        end/.style={ellipse, draw=black, fill=red!30, minimum width=2cm, minimum height=0.8cm, align=center},
        process/.style={rectangle, draw=black, fill=blue!20, minimum width=3cm, minimum height=0.8cm, align=center},
        decision/.style={diamond, draw=black, fill=yellow!30, minimum width=2.5cm, minimum height=2cm, align=center, aspect=2},
        arrow/.style={->, >=stealth, thick}
    ]
    
    % Start
    \node[start] (start) {Эхлэх};
    
    % Process nodes
    \node[process, below=of start] (load) {Flashcard\\ачаалах};
    \node[process, below=of load] (show) {Урд тал\\харуулах};
    \node[decision, below=of show] (flip) {Эргүүлэх\\үү?};
    \node[process, right=3cm of flip] (showback) {Ард тал\\харуулах};
    \node[decision, below=of showback] (rate) {Үнэлгээ\\өгөх};
    \node[process, below=of rate] (update) {Progress\\шинэчлэх};
    \node[decision, below=of update] (more) {Дараагийн\\карт?};
    \node[end, below=of more] (finish) {Дуусгах};
    
    % Arrows
    \draw[arrow] (start) -- (load);
    \draw[arrow] (load) -- (show);
    \draw[arrow] (show) -- (flip);
    \draw[arrow] (flip) -- node[above] {Тийм} (showback);
    \draw[arrow] (flip) -- node[left] {Үгүй} ++(0,-1) -| (show);
    \draw[arrow] (showback) -- (rate);
    \draw[arrow] (rate) -- node[right, font=\tiny] {Амархан/Дунд/Хэцүү} (update);
    \draw[arrow] (update) -- (more);
    \draw[arrow] (more) -- node[right] {Тийм} ++(2,0) |- (load);
    \draw[arrow] (more) -- node[left] {Үгүй} (finish);
    
    \end{tikzpicture}
    \caption{Flashcard суралцах процессын Activity диаграмм}
    \label{fig:activity-flashcard}
\end{figure}

\paragraph{Activity диаграммын тайлбар:}

Диаграмм нь суралцах процессыг харуулна: Flashcard ачаалах → Урд тал харуулах → Эргүүлэх → Ард тал харуулах → Үнэлгээ өгөх (Амархан/Дунд/Хэцүү) → Progress шинэчлэх → Дараагийн карт эсвэл дуусгах. Spaced repetition алгоритм нь хэрэглэгчийн үнэлгээний дагуу дараагийн давталтын огноог тооцоолдог.

\subsubsection{Data Flow Diagram (DFD) - Өгөгдлийн урсгалын диаграмм}

Data Flow Diagram нь системд өгөгдөл хэрхэн урсаж, боловсруулагдаж, хадгалагдаж байгааг харуулдаг. Энэ нь системийн өгөгдлийн архитектурыг ойлгоход тусалдаг.

\begin{figure}[H]
    \centering
    \begin{tikzpicture}[
        node distance=3cm,
        entity/.style={rectangle, draw=black, minimum width=2.5cm, minimum height=1cm, align=center},
        process/.style={circle, draw=black, fill=blue!20, minimum size=2cm, align=center},
        datastore/.style={rectangle, draw=black, fill=green!20, minimum width=2.5cm, minimum height=0.8cm, align=center},
        flow/.style={->, >=stealth, thick}
    ]
    
    % External entities
    \node[entity] (user) {Хэрэглэгч};
    \node[entity, right=5cm of user] (firebase) {Firebase\\Platform};
    
    % Processes
    \node[process, below=of user] (auth) {1.0\\Баталгаа-\\жуулалт};
    \node[process, below=of auth] (study) {2.0\\Суралцах};
    \node[process, right=of study] (ocr) {3.0\\OCR\\таних};
    
    % Data stores
    \node[datastore, below=of study] (userdb) {D1: Users};
    \node[datastore, right=of userdb] (flashdb) {D2: Flashcards};
    \node[datastore, right=of flashdb] (progdb) {D3: UserProgress};
    
    % Data flows
    \draw[flow] (user) -- node[left, font=\tiny] {Email, Password} (auth);
    \draw[flow] (auth) -- node[right, font=\tiny] {User token} (user);
    \draw[flow] (auth) -- node[above, font=\tiny] {Auth request} (firebase);
    \draw[flow] (firebase) -- node[below, font=\tiny] {Token} (auth);
    
    \draw[flow] (user) -- node[left, font=\tiny] {Суралцах хүсэлт} (study);
    \draw[flow] (study) -- node[right, font=\tiny] {Flashcard} (user);
    \draw[flow] (study) -- (flashdb);
    \draw[flow] (flashdb) -- (study);
    \draw[flow] (study) -- (progdb);
    \draw[flow] (progdb) -- (study);
    
    \draw[flow] (user) -- node[above, font=\tiny] {Зураг} (ocr);
    \draw[flow] (ocr) -- node[below, font=\tiny] {Таниулсан үр дүн} (user);
    
    \draw[flow] (auth) -- (userdb);
    \draw[flow] (userdb) -- (auth);
    
    \end{tikzpicture}
    \caption{Egüü системийн Data Flow Diagram (Level 1)}
    \label{fig:dfd-level1}
\end{figure}

\paragraph{Data Flow Diagram-ын тайлбар:}

Диаграмм нь 3 үндсэн процессыг харуулна: (1) Баталгаажуулалт - Firebase-ээр хэрэглэгчийг баталгаажуулах, (2) Суралцах - Flashcards сангаас карт авч UserProgress санд хадгалах, (3) OCR таних - TensorFlow.js ашиглан Монгол бичиг таних. Өгөгдлийн сангууд: Users, Flashcards, UserProgress.



\subsection{Системийн архитектур}
\label{subsec:architecture}

\subsubsection{Системийн ерөнхий архитектур}

Egüü систем нь client-server архитектураар зохион байгуулагдсан бөгөөд дараах үндсэн бүрэлдэхүүн хэсгүүдээс бүрдэнэ:

\begin{figure}[H]
    \centering
    \begin{tikzpicture}[
        node distance=2cm,
        box/.style={rectangle, draw=black, fill=blue!20, minimum width=3cm, minimum height=1cm, align=center},
        db/.style={cylinder, draw=black, fill=green!20, minimum width=2cm, minimum height=1.5cm, align=center, shape border rotate=90},
        cloud/.style={ellipse, draw=black, fill=yellow!20, minimum width=3cm, minimum height=1.5cm, align=center}
    ]
    
    % Client layer
    \node[box] (client) {React Native\\Mobile App};
    
    % Firebase services
    \node[cloud, below=of client] (firebase) {Firebase\\Platform};
    \node[db, below left=1.5cm and 1cm of firebase] (firestore) {Cloud\\Firestore};
    \node[box, below=of firebase] (auth) {Firebase\\Authentication};
    \node[db, below right=1.5cm and 1cm of firebase] (storage) {Firebase\\Storage};
    
    % TensorFlow
    \node[box, right=3cm of client] (tf) {TensorFlow.js\\OCR Model};
    
    % Connections
    \draw[->, thick] (client) -- (firebase);
    \draw[->, thick] (firebase) -- (firestore);
    \draw[->, thick] (firebase) -- (auth);
    \draw[->, thick] (firebase) -- (storage);
    \draw[->, thick] (client) -- (tf);
    
    \end{tikzpicture}
    \caption{Egüü системийн ерөнхий архитектур}
    \label{fig:system-architecture}
\end{figure}

\paragraph{Архитектурын тайлбар:}

\begin{enumerate}
    \item \textbf{Client Layer (React Native Mobile App):}
    \begin{itemize}
        \item iOS болон Android төхөөрөмж дээр ажиллах гар утасны аппликейшн
        \item Expo Router ашиглан навигаци удирдлага
        \item NativeWind/Tailwind CSS ашиглан UI дизайн
        \item React Context API ашиглан төлөв удирдлага
    \end{itemize}
    
    \item \textbf{Firebase Platform:}
    \begin{itemize}
        \item \textbf{Cloud Firestore:} Flashcard, хэрэглэгчийн мэдээлэл, дуртай жагсаалт, streak өгөгдөл хадгалах
        \item \textbf{Firebase Authentication:} Email/Password, Google Sign-In баталгаажуулалт
        \item \textbf{Firebase Storage:} Flashcard зургууд, хэрэглэгчийн профайл зураг хадгалах
    \end{itemize}
    
    \item \textbf{TensorFlow.js OCR Model:}
    \begin{itemize}
        \item Хэрэглэгчийн гараар бичсэн Монгол бичгийг таних
        \item Client-side дээр ажиллах (privacy)
        \item Pre-trained model ашиглах
    \end{itemize}
\end{enumerate}

\subsubsection{Модуляр архитектур}

Egüü апп нь модуляр архитектураар зохион байгуулагдсан бөгөөд дараах үндсэн модулиудаас бүрдэнэ:

\begin{lstlisting}[language=JavaScript, caption=Модуляр бүтэц]
app/
├── (tabs)/                 // Tab navigation
│   ├── index.tsx           // Home screen
│   ├── search.tsx          // Search screen
│   ├── favorites.tsx       // Favorites screen
│   └── profile.tsx         // Profile screen
├── flashcard/
│   └── [id].tsx            // Flashcard detail
├── camera.tsx              // OCR camera screen
├── settings.tsx            // Settings screen
├── _layout.tsx             // Root layout
│
components/                 // Reusable components
├── FlashcardItem.tsx
├── CategoryPill.tsx
├── StreakCounter.tsx
└── DailyWordCard.tsx
│
services/                   // Business logic
├── firebase/
│   ├── auth.ts
│   ├── firestore.ts
│   └── storage.ts
├── ocr/
│   └── tensorflowOCR.ts
└── i18n/
    └── translations.ts
│
contexts/                   // State management
├── AuthContext.tsx
├── FlashcardContext.tsx
└── ThemeContext.tsx
│
utils/                      // Utility functions
├── spacedRepetition.ts
├── dateHelpers.ts
└── validators.ts
\end{lstlisting}

\subsection{Өгөгдлийн сангийн зохиомж}
\label{subsec:database-design}

\subsubsection{Firestore өгөгдлийн бүтэц}

Egüü систем нь Cloud Firestore NoSQL өгөгдлийн санг ашигладаг. Дараах collection-ууд байна:

\paragraph{1. Flashcards Collection:}

\begin{lstlisting}[language=JavaScript, caption=Flashcards collection бүтэц]
flashcards/
  {flashcardId}/
    - mongolBichig: string          // Монгол бичгийн үсэг "ᠠ"
    - phonetic: string               // Фонетик "a"
    - english: string                // Англи орчуулга "a (vowel)"
    - definition: string             // Тодорхойлолт "Эгшиг үсэг"
    - imageUrl: string               // Зургийн URL
    - category: string               // Категори "vowels"
    - difficulty: number             // Хүндрэл (1-5)
    - createdAt: Timestamp
    - updatedAt: Timestamp
\end{lstlisting}

\paragraph{2. Users Collection:}

\begin{lstlisting}[language=JavaScript, caption=Users collection бүтэц]
users/
  {userId}/
    - email: string
    - displayName: string
    - photoURL: string
    - createdAt: Timestamp
    - lastLoginAt: Timestamp
    - preferences:
        - language: string           // "mn" | "en"
        - theme: string              // "light" | "dark"
        - notifications: boolean
    - stats:
        - totalStudied: number       // Нийт суралцсан flashcard
        - currentStreak: number      // Одоогийн streak
        - longestStreak: number      // Хамгийн урт streak
        - lastStudyDate: Timestamp
\end{lstlisting}

\paragraph{3. UserProgress Collection:}

\begin{lstlisting}[language=JavaScript, caption=UserProgress collection бүтэц]
userProgress/
  {userId}/
    flashcards/
      {flashcardId}/
        - easinessFactor: number     // SM-2 алгоритмын EF
        - interval: number           // Дараагийн давталт хүртэлх хоног
        - repetition: number         // Давталтын тоо
        - nextReviewDate: Timestamp  // Дараагийн давталтын огноо
        - lastReviewDate: Timestamp
        - reviewHistory: array       // Давталтын түүх
            - date: Timestamp
            - quality: number        // 0-5
\end{lstlisting}

\paragraph{4. Favorites Collection:}

\begin{lstlisting}[language=JavaScript, caption=Favorites collection бүтэц]
favorites/
  {userId}/
    flashcards: array<string>        // Flashcard ID-ийн жагсаалт
    addedAt: map<string, Timestamp>  // Нэмсэн огноо
\end{lstlisting}

\paragraph{5. DailyWords Collection:}

\begin{lstlisting}[language=JavaScript, caption=DailyWords collection бүтэц]
dailyWords/
  {date}/                            // YYYY-MM-DD формат
    - flashcardId: string
    - createdAt: Timestamp
\end{lstlisting}

\subsubsection{Өгөгдлийн бүтцийн диаграмм}

\begin{figure}[H]
    \centering
    \begin{tikzpicture}[
        node distance=2.5cm,
        entity/.style={rectangle, draw=black, fill=blue!20, minimum width=3cm, minimum height=1.5cm, align=center},
        relationship/.style={diamond, draw=black, fill=yellow!20, minimum width=2cm, minimum height=1cm, align=center}
    ]
    
    % Entities
    \node[entity] (user) {Users};
    \node[entity, right=of user] (flashcard) {Flashcards};
    \node[entity, below=of user] (progress) {UserProgress};
    \node[entity, below=of flashcard] (favorite) {Favorites};
    \node[entity, right=of flashcard] (daily) {DailyWords};
    
    % Relationships
    \draw[->, thick] (user) -- node[above] {has} (progress);
    \draw[->, thick] (user) -- node[above] {has} (favorite);
    \draw[->, thick] (flashcard) -- node[right] {tracked in} (progress);
    \draw[->, thick] (flashcard) -- node[right] {saved in} (favorite);
    \draw[->, thick] (flashcard) -- node[above] {featured in} (daily);
    
    \end{tikzpicture}
    \caption{Өгөгдлийн бүтцийн харилцаа}
    \label{fig:data-model}
\end{figure}

\subsection{Хэрэглэгчийн интерфэйс (UI/UX) загвар}
\label{subsec:ui-ux-design}

\subsubsection{Дизайны зарчмууд}

Egüü апп-ийн UI/UX дизайныг боловсруулахдаа дараах зарчмуудыг баримталсан:

\begin{enumerate}
    \item \textbf{Энгийн, ойлгомжтой:} Шинэ хэрэглэгч хялбар ашиглаж сурах
    \item \textbf{Монгол соёлын элементүүд:} Монгол бичгийн уламжлалт дүрс, өнгө ашиглах
    \item \textbf{Орчин үеийн, сонирхолтой:} Залуу үеийнхний анхаарлыг татах
    \item \textbf{Хүртээмжтэй:} Бүх хэрэглэгчдэд ойлгомжтой (accessibility)
    \item \textbf{Тууштай:} Бүх дэлгэцүүд нэгдмэл дизайнтай
\end{enumerate}

\subsubsection{Өнгөний палитр}

Систем нь орчин үеийн өнгөний схем ашигладаг: Primary Blue (\#3B82F6), Success Green (\#10B981), Error Red (\#EF4444), мөн Light/Dark theme дэмжлэгтэй.

\subsubsection{Үндсэн дэлгэцүүдийн загвар}

Систем нь 4 үндсэн дэлгэцтэй: (1) Нүүр дэлгэц - Өдрийн үг, Streak counter, Категори жагсаалт; (2) Flashcard суралцах дэлгэц - Эргэдэг карт, үнэлгээ өгөх товчнууд; (3) Хайлтын дэлгэц - Хайлтын талбар, категори шүүлтүүр; (4) OCR дэлгэц - Камерын preview, таниулсан үр дүн.

\subsection{Spaced Repetition алгоритмын хэрэгжүүлэлт}
\label{subsec:spaced-repetition-implementation}

Egüü систем нь SuperMemo 2 (SM-2) алгоритмын хялбаршуулсан хувилбарыг ашигладаг. Дараах кодоор хэрэгжүүлсэн:

\begin{lstlisting}[language=JavaScript, caption=SM-2 алгоритмын хэрэгжүүлэлт]
interface CardProgress {
  easinessFactor: number;  // 1.3 - 2.5
  interval: number;        // Хоногоор
  repetition: number;      // Давталтын тоо
  nextReviewDate: Date;
}

function calculateNextReview(
  progress: CardProgress,
  quality: number  // 0-5 (0=буруу, 5=маш амархан)
): CardProgress {
  let { easinessFactor, interval, repetition } = progress;
  
  // Quality 3-аас бага бол эхнээс эхлүүлэх
  if (quality < 3) {
    interval = 1;
    repetition = 0;
  } else {
    // Interval тооцоолох
    if (repetition === 0) {
      interval = 1;
    } else if (repetition === 1) {
      interval = 6;
    } else {
      interval = Math.round(interval * easinessFactor);
    }
    repetition += 1;
  }
  
  // Easiness Factor шинэчлэх
  easinessFactor = easinessFactor + 
    (0.1 - (5 - quality) * (0.08 + (5 - quality) * 0.02));
  
  // EF-ийн доод хязгаар
  if (easinessFactor < 1.3) {
    easinessFactor = 1.3;
  }
  
  // Дараагийн давталтын огноо
  const nextReviewDate = new Date();
  nextReviewDate.setDate(nextReviewDate.getDate() + interval);
  
  return {
    easinessFactor,
    interval,
    repetition,
    nextReviewDate
  };
}
\end{lstlisting}

\subsection{Дүгнэлт}

Энэхүү бүлэгт Egüü системийн шаардлага, архитектур, өгөгдлийн сангийн зохиомж, UI/UX загвар, мөн spaced repetition алгоритмын хэрэгжүүлэлтийг дэлгэрэнгүй танилцуулсан.

Судалгааны арга зүйн үндсэн дүгнэлтүүд:
\begin{itemize}
    \item Системийн шаардлагыг функциональ болон функциональ бус гэж ангилж тодорхойлсон
    \item Client-server архитектураар Firebase платформыг ашиглан backend шийдсэн
    \item NoSQL Firestore өгөгдлийн санг ашиглан уян хатан, scalable өгөгдлийн бүтэц бүтээсэн
    \item Орчин үеийн, хэрэглэхэд хялбар UI/UX дизайн боловсруулсан
    \item SM-2 spaced repetition алгоритмыг хэрэгжүүлж үр дүнтэй сургалтын систем бүтээсэн
\end{itemize}

Дараагийн бүлэгт эдгээр зохиомжийн дагуу системийн бодит хэрэгжүүлэлтийг дэлгэрэнгүй танилцуулна.

\vfill
